\documentclass[aspectratio=169]{beamer}

\usetheme{Madrid}
\usecolortheme{default}

\usepackage{amsmath,amssymb,amsfonts,bm}
\usepackage{mathtools}

\title{The Frobenius Norm and the Trace Identity}
\author{Morten Hjorth-Jensen}
\date{Spring 2026}

\begin{document}

\frame{\titlepage}

\begin{frame}{Goal}
For a real matrix \(A\in\mathbb{R}^{m\times n}\), we want to show that
\[
\|A\|_F^2=\operatorname{Tr}(A^T A),
\qquad
\text{equivalently}
\qquad
\|A\|_F=\sqrt{\operatorname{Tr}(A^T A)}.
\]

We will present two proofs:
\begin{enumerate}
\item a direct component-wise proof,
\item a proof using vectorization and inner products.
\end{enumerate}
\end{frame}

\begin{frame}{Definition of the Frobenius norm}
Let
\[
A=(a_{ij}) \in \mathbb{R}^{m\times n}.
\]

The Frobenius norm is defined by
\[
\|A\|_F
=
\sqrt{\sum_{i=1}^{m}\sum_{j=1}^{n} a_{ij}^2 }.
\]

Therefore,
\[
\|A\|_F^2
=
\sum_{i=1}^{m}\sum_{j=1}^{n} a_{ij}^2.
\]

So the Frobenius norm is simply the Euclidean norm of the matrix entries viewed as one long vector.
\end{frame}

\begin{frame}{Step 1: Entries of \(A^T A\)}
The matrix \(A^T A\) is an \(n\times n\) matrix.

Its \((j,k)\)-entry is
\[
(A^T A)_{jk}
=
\sum_{i=1}^{m} a_{ij}a_{ik}.
\]

In particular, the diagonal entries are
\[
(A^T A)_{jj}
=
\sum_{i=1}^{m} a_{ij}^2.
\]
\end{frame}

\begin{frame}{Step 2: Take the trace}
By definition, the trace is the sum of the diagonal entries:
\[
\operatorname{Tr}(A^T A)
=
\sum_{j=1}^{n}(A^T A)_{jj}.
\]

Using the previous expression,
\[
\operatorname{Tr}(A^T A)
=
\sum_{j=1}^{n}\sum_{i=1}^{m} a_{ij}^2.
\]

Reordering the sums gives
\[
\operatorname{Tr}(A^T A)
=
\sum_{i=1}^{m}\sum_{j=1}^{n} a_{ij}^2.
\]
\end{frame}

\begin{frame}{Conclusion of the direct proof}
But
\[
\sum_{i=1}^{m}\sum_{j=1}^{n} a_{ij}^2
=
\|A\|_F^2.
\]

Hence
\[
\boxed{\|A\|_F^2=\operatorname{Tr}(A^T A)}.
\]

Taking square roots yields
\[
\boxed{\|A\|_F=\sqrt{\operatorname{Tr}(A^T A)}}.
\]
\end{frame}

\begin{frame}{Frobenius inner product}
Define the Frobenius inner product of two matrices \(A,B\in\mathbb{R}^{m\times n}\) by
\[
\langle A,B\rangle_F
=
\operatorname{Tr}(A^T B).
\]

In components,
\[
\operatorname{Tr}(A^T B)
=
\sum_{i=1}^{m}\sum_{j=1}^{n} a_{ij}b_{ij}.
\]

So this is exactly the standard Euclidean inner product of the entries of the two matrices.
\end{frame}

\begin{frame}{The norm induced by the Frobenius inner product}
Every inner product induces a norm:
\[
\|A\|=\sqrt{\langle A,A\rangle}.
\]

For the Frobenius inner product,
\[
\|A\|_F
=
\sqrt{\langle A,A\rangle_F}
=
\sqrt{\operatorname{Tr}(A^T A)}.
\]

Thus the identity
\[
\|A\|_F^2=\operatorname{Tr}(A^T A)
\]
can be viewed as a direct consequence of the fact that the Frobenius norm is the norm induced by the Frobenius inner product.
\end{frame}

\begin{frame}{Vectorization of a matrix}
Define the vectorization map
\[
\operatorname{vec}: \mathbb{R}^{m\times n}\to \mathbb{R}^{mn}
\]
by stacking the columns of \(A\) into one long vector:
\[
\operatorname{vec}(A)
=
\begin{pmatrix}
a_{11}\\
a_{21}\\
\vdots\\
a_{m1}\\
a_{12}\\
a_{22}\\
\vdots\\
a_{mn}
\end{pmatrix}.
\]

The precise ordering is not important for the norm, as long as every matrix entry appears exactly once.
\end{frame}

\begin{frame}{Proof using vectorization}
The Euclidean norm of \(\operatorname{vec}(A)\) is
\[
\|\operatorname{vec}(A)\|_2^2
=
\sum_{i=1}^{m}\sum_{j=1}^{n} a_{ij}^2.
\]

But by definition, this is exactly the Frobenius norm squared:
\[
\|\operatorname{vec}(A)\|_2^2=\|A\|_F^2.
\]

On the other hand, the Euclidean inner product of \(\operatorname{vec}(A)\) with itself is
\[
\operatorname{vec}(A)^T\operatorname{vec}(A).
\]

So
\[
\|A\|_F^2=\operatorname{vec}(A)^T\operatorname{vec}(A).
\]
\end{frame}

\begin{frame}{Vectorization and the trace}
A standard identity is
\[
\operatorname{vec}(A)^T\operatorname{vec}(B)=\operatorname{Tr}(A^T B).
\]

Setting \(B=A\), we get
\[
\operatorname{vec}(A)^T\operatorname{vec}(A)=\operatorname{Tr}(A^T A).
\]

Since
\[
\|A\|_F^2=\operatorname{vec}(A)^T\operatorname{vec}(A),
\]
it follows immediately that
\[
\boxed{\|A\|_F^2=\operatorname{Tr}(A^T A)}.
\]
\end{frame}

\begin{frame}{Why the vectorization identity is true}
Let \(A=(a_{ij})\) and \(B=(b_{ij})\). Then
\[
\operatorname{vec}(A)^T\operatorname{vec}(B)
=
\sum_{i=1}^{m}\sum_{j=1}^{n} a_{ij}b_{ij}.
\]

But from the definition of the trace,
\[
\operatorname{Tr}(A^T B)
=
\sum_{j=1}^{n}(A^T B)_{jj}
=
\sum_{j=1}^{n}\sum_{i=1}^{m} a_{ij}b_{ij}.
\]

Hence
\[
\operatorname{vec}(A)^T\operatorname{vec}(B)=\operatorname{Tr}(A^T B).
\]

So vectorization turns the Frobenius inner product into the standard Euclidean inner product in \(\mathbb{R}^{mn}\).
\end{frame}

\begin{frame}{Geometric interpretation}
The matrix space \(\mathbb{R}^{m\times n}\) is itself a Euclidean vector space.

Under vectorization,
\[
A \longmapsto \operatorname{vec}(A),
\]
the Frobenius norm becomes the ordinary Euclidean norm:
\[
\|A\|_F=\|\operatorname{vec}(A)\|_2.
\]

Thus
\[
\operatorname{Tr}(A^T A)
\]
is nothing but the squared Euclidean length of the matrix when viewed as a vector of all its entries.
\end{frame}

\begin{frame}{Summary}
We have shown that for any real matrix \(A\in\mathbb{R}^{m\times n}\),
\[
\boxed{\|A\|_F^2=\operatorname{Tr}(A^T A)}.
\]

Two proofs were given:
\begin{itemize}
\item a direct component-wise proof using the diagonal entries of \(A^T A\),
\item a conceptual proof using
\[
\operatorname{vec}(A)^T\operatorname{vec}(B)=\operatorname{Tr}(A^T B).
\]
\end{itemize}

Equivalently,
\[
\boxed{\|A\|_F=\sqrt{\operatorname{Tr}(A^T A)}}.
\]
\end{frame}

\end{document}
